
\subsubsection{6.1 Architecture Determinism as the Primary Explanation}

ICC(3,k) = 0.9846 and cross-epsilon tau > 0.95 suggest that the sparsity fingerprint is dominated by the model's weight magnitudes (pre-training geometry) rather than input content. LLaMA-3.1-8B uses SiLU activations (Sigmoid Linear Units) in its MLP blocks, producing soft near-zero values. If the magnitude distribution of gate\_proj weights is the primary determinant of the fraction of activations falling below a fixed threshold, then the sparsity fraction is approximately input-independent by construction. This interpretation is consistent with the observed pattern: the WikiText-103 profiles --- the most distributional-distinct from the other three datasets --- show the lowest tau values relative to the others (0.734--0.750), yet still pass the stability gate comfortably.

The depth gradient --- early layers most sparse, deep layers least sparse --- is consistent with the linguistic hierarchy documented by Clark et al. (2019) and Tenney et al. (2019), in which early layers perform more syntactic, sparse operations and later layers integrate richer semantic representations. This observation is correlational; no causal mechanism is established here.

This architectural determinism interpretation predicts that shuffled or random inputs would produce nearly identical sparsity profiles. That experiment was not conducted; it is noted as a natural validation of the proposed mechanism.

\subsubsection{6.2 Scope of Validated and Unvalidated Claims}

\textbf{Validated (Predictions P1a--P1c):} The fingerprint exists (H-E1: CV = 0.544), is distribution-stable (H-M1: ICC = 0.9846, tau\_min = 0.734), and is threshold-invariant (H-M2: cross-epsilon tau > 0.95). Any reasonable calibration dataset and any epsilon in [0.001, 0.1] yields essentially the same layer rank ordering.

\textbf{Not validated (Predictions P2, P3):} Whether layers with higher sparsity require lower LoRA rank (H-M3, requiring approximately 320 fine-tuning runs: 32 layers $\times$ 5 seeds $\times$ 2 tasks) and whether inverse-sparsity allocation achieves >= 95\% of oracle performance at 60\% parameter budget (H-M4) are empirically open questions. These experiments are fully designed in the research protocol but were not executed in this work. The practical utility of the fingerprint as a rank allocation signal has not been demonstrated.

\subsubsection{6.3 Limitations}

\textbf{L1: Core mechanism unvalidated (severity: high).} The proposed mechanism --- that activation sparsity predicts rank requirements --- has not been tested. Approximately 320 fine-tuning runs are required to test H-M3. The fingerprint's rank-predictive utility is entirely unconfirmed.

\textbf{L2: End-to-end performance unvalidated (severity: high).} No downstream fine-tuning performance result is reported. The claim that inverse-sparsity allocation would achieve efficiency gains at 60\% parameter budget has zero experimental support from this work.

\textbf{L3: SiLU soft-sparsity proxy (severity: medium).} The gate\_proj output is measured before the SiLU activation, so absolute sparsity values represent near-zero pre-activation magnitudes rather than true zero-valued post-activation outputs. Cross-epsilon tau > 0.96 confirms rank ordering robustness to this measurement choice, but absolute sparsity values should not be interpreted as effective sparsity without comparison to effective-rank metrics.

\textbf{L4: Single architecture and task domain (severity: low-medium).} All experiments use LLaMA-3.1-8B. No results are reported for other model families (Mistral, Gemma, Phi). Szatkowski et al. (2025) motivate cross-architecture generalization, but no such generalization is tested here. Similarly, the calibration distributions include GLUE classification tasks but no generation tasks.

\textbf{L5: Protocol deviation --- LLaMA-3.1 vs. LLaMA-3 (severity: low).} The original hypothesis protocol specified LLaMA-3-8B (meta-llama/Meta-Llama-3-8B). LLaMA-3.1-8B was used due to local cache availability. The two models share the same architecture family; this deviation is assessed as minor.

\textbf{L6: Act-LoRA citation unverified (severity: low).} The Act-LoRA citation (actlora2025mdpi) was not verified via Semantic Scholar. Author names, title, volume, and DOI require manual verification before submission.

\subsubsection{6.4 Broader Considerations}

The primary positive implication of this work is that the measurement required is computationally trivial: a single forward pass on 512 samples with no gradient computation, completing in approximately 5 minutes on an NVIDIA H100 NVL using 18.4 GB of VRAM. If H-M3 and H-M4 validate the allocation utility, the zero-cost nature of this measurement makes it broadly applicable to resource-constrained practitioners.

The primary caution is that the current work provides no evidence of improved downstream performance. Practitioners should not use the sparsity fingerprint to inform LoRA rank allocation decisions without first validating the sparsity-rank correlation claimed in H-M3.

